堆和动态极值深度题卡 按题型拆成章内大节，但每个大节都先从 LeetCode 案例切入，再进入分流、案例矩阵、案例推导和实验复盘。这样读者看到的是从具体题推到规律，而不是先读抽象方法再猜它能用在哪。

\topic{堆与动态极值}
这一节先看 LeetCode 23 合并 K 个升序链表、LeetCode 215 数组中的第 K 个最大元素、LeetCode 295 数据流的中位数 这几道 LeetCode 题，再整理同一题型的分流和推导。阅读顺序不是先背原理，而是先看案例的暴力瓶颈，再把重复出现的观察抽象成方法。

\subsection{原理和适用条件}

以 LeetCode 23 合并 K 个升序链表、LeetCode 215 数组中的第 K 个最大元素、LeetCode 295 数据流的中位数 为主线：LeetCode 23 合并 K 个升序链表：模型是“每条链表当前头节点都是可能的下一个最小元素。”，优化抓手是“小顶堆保存每条非空链的当前头，弹出最小节点后把它的后继入堆。”；LeetCode 215 数组中的第 K 个最大元素：模型是“只关心第 K 大，不需要完整排序。”，优化抓手是“大小为 K 的小顶堆保存当前前 K 大边界；快速选择可做到平均线性。”；LeetCode 295 数据流的中位数：模型是“数据流需要动态维护较小一半和较大一半。”，优化抓手是“大顶堆保存较小一半，小顶堆保存较大一半，大小差不超过一。”。这些案例共同推出的规律是：先每轮扫描所有候选找极值，再把动态极值查询交给堆。堆只保证堆顶，不保证全局遍历有序；有过期元素时要在使用堆顶前清理。 方法名只是复盘后的标签，真正要掌握的是从题面约束、暴力失败、状态设计到边界反例的推导链。

\begin{principlebox}{图示：从 LeetCode 案例到 堆与动态极值推导链}
\begin{tabular}{@{}p{0.18\linewidth}p{0.74\linewidth}@{}}
暴力基线 & 枚举候选、完整模拟或逐个检查，先保证覆盖全部可能答案。\\
瓶颈定位 & 慢点在于动态极值查询。我们只需要当前最极端元素，不需要整个集合完全有序。\\
结构选择 & 堆把插入和弹出极值控制在对数复杂度。Top K 用固定大小堆维护边界，合并多路序列用堆保存每一路当前头部，数据流中位数用双堆维护两半。\\
不变量 & 堆不变量不是全局有序，而是堆顶代表当前最优候选。若存在过期元素，需要在弹出时懒删除，保证真正使用堆顶前它仍然有效。\\
代码落地 & C++ 的 \code{std::priority_queue} 默认大顶堆。小顶堆用 \code{std::greater<T>}；结构体比较器要按谁优先级更低来写。堆中保存大对象时尽量保存下标或指针。\\
\end{tabular}
\end{principlebox}

\subsection{LeetCode 案例分流}

\begin{principlebox}{从 LeetCode 案例分流：堆与动态极值}
\textbf{Top K 和动态排名。}
代表案例：LeetCode 215 数组第 K 大、LeetCode 347 前 K 高频、LeetCode 973 最接近原点 K 个点。先看这些题的题面条件：只需要前 K 个、最近 K 个或第 K 个，不需要完整排序。由此推出的处理方式是：维护固定大小堆，堆顶是当前边界。风险点：K 的边界、比较器方向和输出顺序要说明。

\textbf{多路归并。}
代表案例：LeetCode 23 合并 K 个升序链表、LeetCode 632 最小区间覆盖 K 个列表。先看这些题的题面条件：多个有序序列各自只能暴露当前头部。由此推出的处理方式是：堆保存每一路当前候选，弹出后推进该路。风险点：堆里最多每路一个候选，某一路耗尽会改变终止条件。

\textbf{双堆和数据流。}
代表案例：LeetCode 295 数据流中位数、LeetCode 480 滑动窗口中位数。先看这些题的题面条件：数据在线到达，需要动态维护中位数或两侧平衡。由此推出的处理方式是：大顶堆保存较小一半，小顶堆保存较大一半。风险点：若支持删除，需要延迟删除，堆顶使用前必须清理。

\textbf{调度和反悔。}
代表案例：LeetCode 621 任务调度器、LeetCode 630 课程表 III、LeetCode 1675 数组最小偏移量。先看这些题的题面条件：每一步选择当前收益最大、冷却结束或最紧急候选。由此推出的处理方式是：用堆组织当前可选任务或可反悔候选。风险点：要区分堆模拟和数学公式，二者适用条件不同。

\end{principlebox}

\subsection{案例矩阵}

本节案例覆盖 LeetCode 23 合并 K 个升序链表、LeetCode 215 数组中的第 K 个最大元素、LeetCode 295 数据流的中位数、LeetCode 347 前 K 个高频元素、LeetCode 480 滑动窗口中位数、LeetCode 502 IPO、LeetCode 621 任务调度器、LeetCode 632 最小区间覆盖 K 个列表、LeetCode 703 数据流中的第 K 大元素、LeetCode 767 重构字符串、LeetCode 857 雇佣 K 名工人的最低成本、LeetCode 871 最低加油次数、LeetCode 973 最接近原点的 K 个点、LeetCode 1046 最后一块石头的重量、LeetCode 1353 最多可以参加的会议数目、LeetCode 1675 数组的最小偏移量、LeetCode 1851 包含每个查询的最小区间。阅读时不要按题号背答案，而要先判断题面落在哪个分流场景：查询目标是什么，输入约束给了什么单调性或等价关系，暴力慢在哪里，优化结构保存了什么状态。

\begin{principlebox}{案例矩阵}
\textbf{LeetCode 23 合并 K 个升序链表。}
每条链表当前头节点都是可能的下一个最小元素。单点风险：空链表数组、某条链为空、重复值、堆里保存节点指针还是值。

\textbf{LeetCode 215 数组中的第 K 个最大元素。}
只关心第 K 大，不需要完整排序。单点风险：K 为一、K 为 n、重复值、随机 pivot。

\textbf{LeetCode 295 数据流的中位数。}
数据流需要动态维护较小一半和较大一半。单点风险：偶数个元素、负数、重复值、平均值溢出。

\textbf{LeetCode 347 前 K 个高频元素。}
先统计频次，再只维护前 K 个候选。单点风险：K 等于不同元素数、频次相同、结果顺序不要求。

\textbf{LeetCode 480 滑动窗口中位数。}
窗口移动时既要加入新元素，也要删除滑出的旧元素并查询中位数。单点风险：重复值、偶数窗口、平均值溢出、堆顶过期。

\textbf{LeetCode 502 IPO。}
当前资本决定可解锁项目集合，目标是在最多 k 次选择中最大化最终资本。单点风险：初始资本为零、没有可做项目、k 大于项目数、利润为零、资本解锁链式增长。

\textbf{LeetCode 621 任务调度器。}
相同任务之间有冷却间隔，核心受最高频任务骨架限制。单点风险：冷却为零、多个最高频、空闲被其他任务填满。

\textbf{LeetCode 632 最小区间覆盖 K 个列表。}
每个列表至少取一个元素时，当前区间由最小当前值和最大当前值决定。单点风险：某个列表为空、重复值、区间长度相同的 tie-breaker、列表耗尽。

\textbf{LeetCode 703 数据流中的第 K 大元素。}
数据不断到达，只关心当前第 K 大而非完整排序。单点风险：初始元素少于 K、重复值、K 为一、连续 add。

\textbf{LeetCode 767 重构字符串。}
相邻不能相同，剩余次数最多的字符最难安排。单点风险：单字符、最高频过半、多个字符同频、最后只剩一个字符、结果不唯一。

\textbf{LeetCode 857 雇佣 K 名工人的最低成本。}
工资比例由当前队伍中最高工资质量比决定，总成本等于比例乘质量和。单点风险：K 为一、比例相同、浮点精度、质量和更新。

\textbf{LeetCode 871 最低加油次数。}
行驶过程中，经过的加油站形成可反悔的燃料候选。单点风险：起点油足够、无法到达第一个站、多个站同位置、目标作为终点事件。

\textbf{LeetCode 973 最接近原点的 K 个点。}
距离只需比较平方，避免开方误差和成本。单点风险：K 为一、K 为全部、距离相同、坐标平方溢出。

\textbf{LeetCode 1046 最后一块石头的重量。}
每次都要取当前最重的两块石头碰撞。单点风险：没有石头、一块石头、两块相等、多次产生相同重量。

\textbf{LeetCode 1353 最多可以参加的会议数目。}
每天只能参加一个正在进行的会议，应优先选择最早结束的会议。单点风险：同一天多个会议、过期会议清理、空闲日期跳跃、结束日相同。

\textbf{LeetCode 1675 数组的最小偏移量。}
奇数可以先乘二，偶数可以不断除二，目标是缩小当前最大最小差。单点风险：全奇数、全偶数、最大值变奇数停止、数值范围。

\textbf{LeetCode 1851 包含每个查询的最小区间。}
每个查询点需要找覆盖它的最短区间，区间和查询都可离线排序。单点风险：没有覆盖区间、区间长度相同、查询重复、右端过期。

\end{principlebox}

\subsection{共性落点}

\begin{principlebox}{本组共性落点}
\textbf{慢版本。} 暴力做法每轮扫描所有候选找最大或最小，或者每次变化后重新排序。若候选集合动态变化，这会把大量旧排序工作丢掉。
\textbf{瓶颈。} 慢点在于动态极值查询。我们只需要当前最极端元素，不需要整个集合完全有序。
\textbf{状态字段。} 堆元素字段、堆顶比较规则、候选是否过期、答案读取时机；如果需要懒删除，还要保存删除计数。
\textbf{不变量。} 堆不变量不是全局有序，而是堆顶代表当前最优候选。若存在过期元素，需要在弹出时懒删除，保证真正使用堆顶前它仍然有效。
\textbf{实现检查。} 比较器方向先用小样例验证；每次使用堆顶前清理过期项；堆中保存下标时要能回查原数据。
\textbf{C++ 落地。} C++ 的 \code{std::priority_queue} 默认大顶堆。小顶堆用 \code{std::greater<T>}；结构体比较器要按谁优先级更低来写。堆中保存大对象时尽量保存下标或指针。
\textbf{证明抓手。} 证明从候选覆盖开始：所有可能成为下一步答案的对象都在堆中，堆顶过期时不会被使用。
\end{principlebox}

\subsection{案例推导}

\noindent\textbf{LeetCode 23 合并 K 个升序链表。}

\textbf{题目说明。} LeetCode 23 合并 K 个升序链表 的题面入口要先还原成具体任务：每条链表当前头节点都是可能的下一个最小元素。

\textbf{样例入口。} 拿三条链表头 1、1、2 手算，每次答案只需要当前最小头节点；弹出某条链的头后，只有它的后继会成为新候选。暴力每轮扫 k 个头，堆把“找当前最小头”变成对数操作。

\textbf{暴力基线。} 暴力每轮扫描 K 个链表头，找到当前最小头节点接到结果后推进该链。

\textbf{暴力过程。} 若三条链当前头是 1、1、2，第一步只需从这三个头里选最小；弹出某条链的 1 后，只有它的后继会成为新候选。

\textbf{重复工作。} 题面模型：每条链表当前头节点都是可能的下一个最小元素。暴力版的慢点要落到本题：浪费在每输出一个节点都线性扫 K 个头。候选集合动态变化，但每次只需要最小头节点，适合小顶堆。后面的优化只消掉这类重复工作，不能改变答案集合。

\textbf{优化条件。} 本题的关键观察是：小顶堆保存每条非空链的当前头，弹出最小节点后把它的后继入堆。这句话必须能翻译成可维护状态；如果题目条件改变后观察不再成立，就不能继续套原来的写法。

\textbf{相邻状态。} 规律是堆里始终最多保存每条非空链的一个当前头。把这个特例继续拆开看：堆顶必须正好代表下一步最需要处理的候选；若堆里可能有过期对象，读取堆顶前要先清理。堆只解决动态极值，不自动证明被弹出或被丢弃的元素无用。

\textbf{状态复用。} 小顶堆保存每条非空链的当前头节点。弹出最小节点接入结果；若该节点有 next，把 next 入堆。手算时只改被新输入影响的那一小部分，而不是回头重算完整候选。

\textbf{指针和字段。} 状态字段要能说清：heap 中元素是 ListNode*，比较依据是 node->val；tail 是结果尾；node->next 是该链的新候选。手算时按这个协议跟踪：头节点 1、1、2 入堆；弹出第一个 1 后把它的后继入堆，堆仍保存每条链当前最小暴露节点。

\textbf{代码落点。} 关键是堆里保存节点指针而不是复制值，这样可以复用原节点。比较器要写成小顶堆语义，C++ priority\_queue 默认是大顶堆。关键代码行必须能逐句翻译成“旧状态怎样变成新状态”，否则就只是模板。

\textbf{正确性。} 每条有序链未输出部分的最小值就是当前头；所有链当前头中的最小者就是全局下一个输出节点。

\textbf{边界实验。} 先用空链表数组、某条链为空、重复值、堆里保存节点指针还是值做反例检查。实验时覆盖空链表数组、某条链为空、重复值、堆里保存节点指针还是值，记录堆顶是否过期、比较器是否让正确候选在堆顶、K 或窗口大小为边界值时是否稳定。

\textbf{复杂度变化。} 每轮扫 K 个头是 O(NK)。堆中最多 K 个节点，每个节点入堆出堆一次，总时间 O(N log K)，空间 O(K)。

\textbf{迁移追问。} 如果题目改成“也可以分治两两合并，堆更适合 K 较大且链长不均的场景”，先判断原状态是否仍然覆盖新答案集合，再决定是微调边界、改 value 语义，还是换算法。

\noindent\textbf{LeetCode 215 数组中的第 K 个最大元素。}

\textbf{题目说明。} LeetCode 215 数组中的第 K 个最大元素 的题面入口要先还原成具体任务：只关心第 K 大，不需要完整排序。

\textbf{样例入口。} 拿 [3, 2, 1, 5, 6, 4]、k=2 手算，固定大小为 2 的小顶堆最终保存最大的两个数，堆顶就是第 2 大。

\textbf{暴力基线。} 暴力排序整个数组，然后返回降序第 k 个或升序第 n-k 个。

\textbf{暴力过程。} 以 [3, 2, 1, 5, 6, 4]、k=2 为例，完整排序后第 2 大是 5；但为了一个第 K 大，不需要知道所有元素的完整顺序。

\textbf{重复工作。} 题面模型：只关心第 K 大，不需要完整排序。暴力版的慢点要落到本题：浪费在完整排序。我们只需要维护当前最大的 K 个元素，其中最小的那个就是第 K 大边界。后面的优化只消掉这类重复工作，不能改变答案集合。

\textbf{优化条件。} 本题的关键观察是：大小为 K 的小顶堆保存当前前 K 大边界；快速选择可做到平均线性。这句话必须能翻译成可维护状态；如果题目条件改变后观察不再成立，就不能继续套原来的写法。

\textbf{相邻状态。} 规律是堆里只保留当前前 k 大候选，比堆顶小的数不可能进入答案。把这个特例继续拆开看：堆顶必须正好代表下一步最需要处理的候选；若堆里可能有过期对象，读取堆顶前要先清理。堆只解决动态极值，不自动证明被弹出或被丢弃的元素无用。

\textbf{状态复用。} 大小为 K 的小顶堆保存当前前 K 大。新数大于堆顶时，弹出堆顶并加入新数；小于等于堆顶则不可能进入前 K。手算时只改被新输入影响的那一小部分，而不是回头重算完整候选。

\textbf{指针和字段。} 状态字段要能说清：heap.top 是当前前 K 大中的最小值，size 保持不超过 K，num 是当前扫描元素。手算时按这个协议跟踪：扫描样例时，堆最终保存 5 和 6，堆顶 5 就是第 2 大。

\textbf{代码落点。} 关键是使用小顶堆，C++ priority\_queue 默认大顶堆，需要 greater<int>。若选择快速选择，要处理随机 pivot 和重复值。关键代码行必须能逐句翻译成“旧状态怎样变成新状态”，否则就只是模板。

\textbf{正确性。} 堆中始终保存已扫描元素的前 K 大；若新元素不超过堆顶，它最多排在第 K 大之后，可以安全丢弃。

\textbf{边界实验。} 先用K 为一、K 为 n、重复值、随机 pivot做反例检查。实验时覆盖K 为一、K 为 n、重复值、随机 pivot，记录堆顶是否过期、比较器是否让正确候选在堆顶、K 或窗口大小为边界值时是否稳定。

\textbf{复杂度变化。} 排序 O(n log n)，小顶堆 O(n log k) 时间、O(k) 空间；快速选择平均 O(n)。

\textbf{迁移追问。} 如果题目改成“若数据流持续到来，堆方案比一次性快速选择更自然”，先判断原状态是否仍然覆盖新答案集合，再决定是微调边界、改 value 语义，还是换算法。

\noindent\textbf{LeetCode 295 数据流的中位数。}

\textbf{题目说明。} LeetCode 295 数据流的中位数 的题面入口要先还原成具体任务：数据流需要动态维护较小一半和较大一半。

\textbf{样例入口。} 拿数据流 1、2、3 手算，插入 1 后中位数 1；插入 2 后两半是 [1] 和 [2]，中位数 1.5；插入 3 后右半多一个，要平衡成左半 [2, 1]、右半 [3]。

\textbf{暴力基线。} 暴力每插入一个数就重新排序所有已到达元素，然后按长度奇偶取中位数。

\textbf{暴力过程。} 数据流 1、2、3 中，插入 1 中位数是 1；插入 2 后中位数是 (1+2)/2；插入 3 后中位数是 2。

\textbf{重复工作。} 题面模型：数据流需要动态维护较小一半和较大一半。暴力版的慢点要落到本题：浪费在每次都维护完整排序。中位数只需要较小一半的最大值和较大一半的最小值。后面的优化只消掉这类重复工作，不能改变答案集合。

\textbf{优化条件。} 本题的关键观察是：大顶堆保存较小一半，小顶堆保存较大一半，大小差不超过一。这句话必须能翻译成可维护状态；如果题目条件改变后观察不再成立，就不能继续套原来的写法。

\textbf{相邻状态。} 规律是大顶堆保存较小一半，小顶堆保存较大一半，大小差最多 1。把这个特例继续拆开看：堆顶必须正好代表下一步最需要处理的候选；若堆里可能有过期对象，读取堆顶前要先清理。堆只解决动态极值，不自动证明被弹出或被丢弃的元素无用。

\textbf{状态复用。} 大顶堆 small 保存较小一半，小顶堆 large 保存较大一半。保持 small.size() 与 large.size() 差不超过 1，且 small.top()<=large.top()。手算时只改被新输入影响的那一小部分，而不是回头重算完整候选。

\textbf{指针和字段。} 状态字段要能说清：small.top 是左半最大值，large.top 是右半最小值；奇数时较大堆顶是中位数，偶数时两个堆顶平均。手算时按这个协议跟踪：插入 1 放 small；插入 2 后移动平衡成 small=[1]、large=[2]；插入 3 后 large 多，再移动 2 到 small，中位数 2。

\textbf{代码落点。} 关键是先按顺序不变量放入堆，再重平衡大小。求平均时用 double，并防止两个 int 相加溢出。关键代码行必须能逐句翻译成“旧状态怎样变成新状态”，否则就只是模板。

\textbf{正确性。} 两个堆按大小分割数据流，中位数只可能位于左半最大或右半最小；大小不变量保证取堆顶即可。

\textbf{边界实验。} 先用偶数个元素、负数、重复值、平均值溢出做反例检查。实验时覆盖偶数个元素、负数、重复值、平均值溢出，记录堆顶是否过期、比较器是否让正确候选在堆顶、K 或窗口大小为边界值时是否稳定。

\textbf{复杂度变化。} 每次重新排序 O(n log n)；双堆 addNum O(log n)，findMedian O(1)，空间 O(n)。

\textbf{迁移追问。} 如果题目改成“若还要删除任意元素，需要延迟删除哈希表”，先判断原状态是否仍然覆盖新答案集合，再决定是微调边界、改 value 语义，还是换算法。

\noindent\textbf{LeetCode 347 前 K 个高频元素。}

\textbf{题目说明。} LeetCode 347 前 K 个高频元素 的题面入口要先还原成具体任务：先统计频次，再只维护前 K 个候选。

\textbf{样例入口。} 拿 [1, 1, 1, 2, 2, 3]、k=2 手算，频次是 1->3、2->2、3->1，前两个高频是 1 和 2。

\textbf{暴力基线。} 暴力先统计频次，再把所有不同元素按频次完整排序，取前 K 个。

\textbf{暴力过程。} 以 [1, 1, 1, 2, 2, 3]、k=2 为例，频次为 1->3、2->2、3->1，答案是 1 和 2。

\textbf{重复工作。} 题面模型：先统计频次，再只维护前 K 个候选。暴力版的慢点要落到本题：浪费在完整排序所有不同元素。只要前 K 高频，低于当前第 K 高频边界的元素可以丢弃。后面的优化只消掉这类重复工作，不能改变答案集合。

\textbf{优化条件。} 本题的关键观察是：小顶堆大小超过 K 就弹出最低频；桶排序利用频次上界。这句话必须能翻译成可维护状态；如果题目条件改变后观察不再成立，就不能继续套原来的写法。

\textbf{相邻状态。} 规律是先用哈希计数，再用大小为 k 的小顶堆或桶按频次取 Top K；频次相同不影响题目集合答案。把这个特例继续拆开看：堆顶必须正好代表下一步最需要处理的候选；若堆里可能有过期对象，读取堆顶前要先清理。堆只解决动态极值，不自动证明被弹出或被丢弃的元素无用。

\textbf{状态复用。} 先用哈希表计数。再用大小为 K 的小顶堆保存候选，堆顶是当前前 K 中频次最低者；超过 K 就弹出。手算时只改被新输入影响的那一小部分，而不是回头重算完整候选。

\textbf{指针和字段。} 状态字段要能说清：freq[value] 是频次，heap 元素是 (frequency, value)，heap.top 是淘汰边界。手算时按这个协议跟踪：插入频次 3 和 2 后堆满；频次 1 的元素入堆会被立刻弹出，留下 1 和 2。

\textbf{代码落点。} 关键是比较器按频次建小顶堆。题目通常不要求输出顺序；若要求顺序，最后还要按规则排序结果。关键代码行必须能逐句翻译成“旧状态怎样变成新状态”，否则就只是模板。

\textbf{正确性。} 堆始终保存已处理元素中频次最高的 K 个；任何被弹出的元素频次不高于堆中 K 个候选，不可能进入答案。

\textbf{边界实验。} 先用K 等于不同元素数、频次相同、结果顺序不要求做反例检查。实验时覆盖K 等于不同元素数、频次相同、结果顺序不要求，记录堆顶是否过期、比较器是否让正确候选在堆顶、K 或窗口大小为边界值时是否稳定。

\textbf{复杂度变化。} 完整排序 O(m log m)，m 为不同元素数；堆法 O(n+m log k)，空间 O(m+k)。

\textbf{迁移追问。} 如果题目改成“若要按频次和字典序排序，比较器要同时处理两个键”，先判断原状态是否仍然覆盖新答案集合，再决定是微调边界、改 value 语义，还是换算法。

\noindent\textbf{LeetCode 480 滑动窗口中位数。}

\textbf{题目说明。} LeetCode 480 滑动窗口中位数 的题面入口要先还原成具体任务：窗口移动时既要加入新元素，也要删除滑出的旧元素并查询中位数。

\textbf{样例入口。} 拿 [1, 3, -1]、k=3 手算，中位数是 1；窗口右移时要删除 1、加入 -3，普通堆不能任意删除。

\textbf{暴力基线。} 慢版本每轮扫描全部候选来找最大或最小，再更新候选集合。它的答案选择过程清楚，但动态集合每变化一次就重新找极值，代价太高。放到本题就是：先按“窗口移动时既要加入新元素，也要删除滑出的旧元素并查询中位数”完整试慢版本；样例里已经能看到“规律是双堆维护两半，再用延迟删除清理过期元素，堆顶使用前必须有效”，所以优化前必须先能说清慢版本枚举了哪些候选。

\textbf{暴力过程。} 拿 [1, 3, -1]、k=3 手算，中位数是 1；窗口右移时要删除 1、加入 -3，普通堆不能任意删除。

\textbf{重复工作。} 题面模型：窗口移动时既要加入新元素，也要删除滑出的旧元素并查询中位数。暴力版的慢点要落到本题：慢版本会围绕“窗口移动时既要加入新元素，也要删除滑出的旧元素并查询中位数”反复扫描动态候选集找极值。样例规律是“规律是双堆维护两半，再用延迟删除清理过期元素，堆顶使用前必须有效”，说明每轮真正需要的是堆顶边界，而不是把候选全集重新排序。后面的优化只消掉这类重复工作，不能改变答案集合。

\textbf{优化条件。} 本题的关键观察是：双堆维护左右两半，延迟删除表记录已离窗但尚未到堆顶的元素。这句话必须能翻译成可维护状态；如果题目条件改变后观察不再成立，就不能继续套原来的写法。

\textbf{相邻状态。} 规律是双堆维护两半，再用延迟删除清理过期元素，堆顶使用前必须有效。把这个特例继续拆开看：堆顶必须正好代表下一步最需要处理的候选；若堆里可能有过期对象，读取堆顶前要先清理。堆只解决动态极值，不自动证明被弹出或被丢弃的元素无用。

\textbf{状态复用。} 把每轮全量扫描改成堆顶查询；插入、弹出和过期清理共同维护动态候选集。本题落点是：规律是双堆维护两半，再用延迟删除清理过期元素，堆顶使用前必须有效。手算时只改被新输入影响的那一小部分，而不是回头重算完整候选。

\textbf{指针和字段。} 状态字段要能说清：堆元素字段、堆顶比较规则、候选是否过期、答案读取时机；如果需要懒删除，还要保存删除计数。手算时按这个协议跟踪：拿 [1, 3, -1]、k=3 手算，中位数是 1；窗口右移时要删除 1、加入 -3，普通堆不能任意删除。然后按协议复查：手算时记录堆顶、候选集合和是否有过期元素。每次使用堆顶前都要确认它仍然属于当前问题。

\textbf{代码落点。} 比较器方向先用小样例验证；每次使用堆顶前清理过期项；堆中保存下标时要能回查原数据。关键代码行必须能逐句翻译成“旧状态怎样变成新状态”，否则就只是模板。

\textbf{正确性。} 证明从候选覆盖开始：所有可能成为下一步答案的对象都在堆中，堆顶过期时不会被使用。

\textbf{边界实验。} 先用重复值、偶数窗口、平均值溢出、堆顶过期做反例检查。实验时覆盖重复值、偶数窗口、平均值溢出、堆顶过期，记录堆顶是否过期、比较器是否让正确候选在堆顶、K 或窗口大小为边界值时是否稳定。

\textbf{复杂度变化。} 暴力每轮扫描或排序动态候选，会把线性扫描或排序反复发生；堆把取极值、插入和删除边界控制在对数级。

\textbf{迁移追问。} 如果题目改成“普通 priority\_queue 不能任意删除，所以必须解释懒删除”，先判断原状态是否仍然覆盖新答案集合，再决定是微调边界、改 value 语义，还是换算法。

\noindent\textbf{LeetCode 502 IPO。}

\textbf{题目说明。} LeetCode 502 IPO 的题面入口要先还原成具体任务：当前资本决定可解锁项目集合，目标是在最多 k 次选择中最大化最终资本。

\textbf{样例入口。} 拿 k=2、w=0、profits=[1, 2, 3]、capital=[0, 1, 1] 手算，初始只能做资本 0 的项目，做完利润 1 后资本变 1，才能解锁利润 2 和 3 的项目。

\textbf{暴力基线。} 慢版本每轮扫描全部候选来找最大或最小，再更新候选集合。它的答案选择过程清楚，但动态集合每变化一次就重新找极值，代价太高。放到本题就是：先按“当前资本决定可解锁项目集合，目标是在最多 k 次选择中最大化最终资本”完整试慢版本；样例里已经能看到“规律是按资本排序解锁候选，用利润大顶堆每轮选择当前可做项目里的最大利润”，所以优化前必须先能说清慢版本枚举了哪些候选。

\textbf{暴力过程。} 拿 k=2、w=0、profits=[1, 2, 3]、capital=[0, 1, 1] 手算，初始只能做资本 0 的项目，做完利润 1 后资本变 1，才能解锁利润 2 和 3 的项目。

\textbf{重复工作。} 题面模型：当前资本决定可解锁项目集合，目标是在最多 k 次选择中最大化最终资本。暴力版的慢点要落到本题：慢版本会围绕“当前资本决定可解锁项目集合，目标是在最多 k 次选择中最大化最终资本”反复扫描动态候选集找极值。样例规律是“规律是按资本排序解锁候选，用利润大顶堆每轮选择当前可做项目里的最大利润”，说明每轮真正需要的是堆顶边界，而不是把候选全集重新排序。后面的优化只消掉这类重复工作，不能改变答案集合。

\textbf{优化条件。} 本题的关键观察是：项目按 capital 排序，用指针把可负担项目加入利润大顶堆，每轮选利润最高。这句话必须能翻译成可维护状态；如果题目条件改变后观察不再成立，就不能继续套原来的写法。

\textbf{相邻状态。} 规律是按资本排序解锁候选，用利润大顶堆每轮选择当前可做项目里的最大利润。把这个特例继续拆开看：堆顶必须正好代表下一步最需要处理的候选；若堆里可能有过期对象，读取堆顶前要先清理。堆只解决动态极值，不自动证明被弹出或被丢弃的元素无用。

\textbf{状态复用。} 把每轮全量扫描改成堆顶查询；插入、弹出和过期清理共同维护动态候选集。本题落点是：规律是按资本排序解锁候选，用利润大顶堆每轮选择当前可做项目里的最大利润。手算时只改被新输入影响的那一小部分，而不是回头重算完整候选。

\textbf{指针和字段。} 状态字段要能说清：堆元素字段、堆顶比较规则、候选是否过期、答案读取时机；如果需要懒删除，还要保存删除计数。手算时按这个协议跟踪：拿 k=2、w=0、profits=[1, 2, 3]、capital=[0, 1, 1] 手算，初始只能做资本 0 的项目，做完利润 1 后资本变 1，才能解锁利润 2 和 3 的项目。然后按协议复查：手算时记录堆顶、候选集合和是否有过期元素。每次使用堆顶前都要确认它仍然属于当前问题。

\textbf{代码落点。} 比较器方向先用小样例验证；每次使用堆顶前清理过期项；堆中保存下标时要能回查原数据。关键代码行必须能逐句翻译成“旧状态怎样变成新状态”，否则就只是模板。

\textbf{正确性。} 证明从候选覆盖开始：所有可能成为下一步答案的对象都在堆中，堆顶过期时不会被使用。

\textbf{边界实验。} 先用初始资本为零、没有可做项目、k 大于项目数、利润为零、资本解锁链式增长做反例检查。实验时覆盖初始资本为零、没有可做项目、k 大于项目数、利润为零、资本解锁链式增长，记录堆顶是否过期、比较器是否让正确候选在堆顶、K 或窗口大小为边界值时是否稳定。

\textbf{复杂度变化。} 暴力每轮扫描或排序动态候选，会把线性扫描或排序反复发生；堆把取极值、插入和删除边界控制在对数级。

\textbf{迁移追问。} 如果题目改成“若项目有依赖关系，需要拓扑或约束优化，不能只按资本解锁”，先判断原状态是否仍然覆盖新答案集合，再决定是微调边界、改 value 语义，还是换算法。

\noindent\textbf{LeetCode 621 任务调度器。}

\textbf{题目说明。} LeetCode 621 任务调度器 的题面入口要先还原成具体任务：相同任务之间有冷却间隔，核心受最高频任务骨架限制。

\textbf{样例入口。} 拿 A, A, A, B, B, B、n=2 手算，若只按顺序执行 A 会违反冷却，必须在相同任务之间间隔两个时间单位。

\textbf{暴力基线。} 暴力按时间模拟，每个时刻扫描所有任务，找一个未冷却且剩余次数最多的任务执行，否则空闲。

\textbf{暴力过程。} 任务 A, A, A, B, B, B、n=2 中，相同 A 之间必须间隔两个单位；一种最短安排是 A B idle A B idle A B。

\textbf{重复工作。} 题面模型：相同任务之间有冷却间隔，核心受最高频任务骨架限制。暴力版的慢点要落到本题：浪费在逐时刻扫描所有任务。更重要的是最高频任务决定冷却骨架，其他任务只是填空。后面的优化只消掉这类重复工作，不能改变答案集合。

\textbf{优化条件。} 本题的关键观察是：可用大顶堆模拟，也可用最高频公式计算最短时间。这句话必须能翻译成可维护状态；如果题目条件改变后观察不再成立，就不能继续套原来的写法。

\textbf{相邻状态。} 规律是最高频任务决定骨架长度，也可用堆按剩余次数模拟，每轮安排最多 n+1 个不同任务。把这个特例继续拆开看：堆顶必须正好代表下一步最需要处理的候选；若堆里可能有过期对象，读取堆顶前要先清理。堆只解决动态极值，不自动证明被弹出或被丢弃的元素无用。

\textbf{状态复用。} 公式法统计最大频次 maxFreq 和拥有最大频次的任务数 maxCount。骨架长度为 (maxFreq-1)*(n+1)+maxCount，答案取它和任务总数的较大值。手算时只改被新输入影响的那一小部分，而不是回头重算完整候选。

\textbf{指针和字段。} 状态字段要能说清：maxFreq 是最高频任务次数，maxCount 是最高频任务种类数，n 是冷却间隔，tasks.size() 是没有空闲时的下界。手算时按这个协议跟踪：A 和 B 都出现 3 次，maxFreq=3、maxCount=2，骨架为 (3-1)*3+2=8。

\textbf{代码落点。} 关键是最后 return max(tasks.size(), skeleton)。若 n=0，任务总数就是答案。关键代码行必须能逐句翻译成“旧状态怎样变成新状态”，否则就只是模板。

\textbf{正确性。} 最高频任务之间至少有 maxFreq-1 个间隔块，每块长度 n+1；多个最高频任务占据最后一列。其他任务只能填充空位，不能突破任务总数下界。

\textbf{边界实验。} 先用冷却为零、多个最高频、空闲被其他任务填满做反例检查。实验时覆盖冷却为零、多个最高频、空闲被其他任务填满，记录堆顶是否过期、比较器是否让正确候选在堆顶、K 或窗口大小为边界值时是否稳定。

\textbf{复杂度变化。} 堆模拟 O(T log C)，公式法统计频次 O(T+C)，空间 O(C)，C 为任务种类数。

\textbf{迁移追问。} 如果题目改成“若任务有不同释放时间或执行时长，就变成更一般的调度问题”，先判断原状态是否仍然覆盖新答案集合，再决定是微调边界、改 value 语义，还是换算法。

\noindent\textbf{LeetCode 632 最小区间覆盖 K 个列表。}

\textbf{题目说明。} LeetCode 632 最小区间覆盖 K 个列表 的题面入口要先还原成具体任务：每个列表至少取一个元素时，当前区间由最小当前值和最大当前值决定。

\textbf{样例入口。} 拿三组 [4, 10]、[0, 9]、[5, 18] 手算，当前头是 4、0、5，区间 [0, 5] 覆盖三组；弹出最小头 0 后推进该组。

\textbf{暴力基线。} 慢版本每轮扫描全部候选来找最大或最小，再更新候选集合。它的答案选择过程清楚，但动态集合每变化一次就重新找极值，代价太高。放到本题就是：先按“每个列表至少取一个元素时，当前区间由最小当前值和最大当前值决定”完整试慢版本；样例里已经能看到“规律是堆保存每组当前元素，同时维护当前最大值，最小头决定下一步推进哪一组”，所以优化前必须先能说清慢版本枚举了哪些候选。

\textbf{暴力过程。} 拿三组 [4, 10]、[0, 9]、[5, 18] 手算，当前头是 4、0、5，区间 [0, 5] 覆盖三组；弹出最小头 0 后推进该组。

\textbf{重复工作。} 题面模型：每个列表至少取一个元素时，当前区间由最小当前值和最大当前值决定。暴力版的慢点要落到本题：慢版本会围绕“每个列表至少取一个元素时，当前区间由最小当前值和最大当前值决定”反复扫描动态候选集找极值。样例规律是“规律是堆保存每组当前元素，同时维护当前最大值，最小头决定下一步推进哪一组”，说明每轮真正需要的是堆顶边界，而不是把候选全集重新排序。后面的优化只消掉这类重复工作，不能改变答案集合。

\textbf{优化条件。} 本题的关键观察是：小顶堆保存各列表当前元素，同时维护当前最大值；弹出最小所在列表并推进。这句话必须能翻译成可维护状态；如果题目条件改变后观察不再成立，就不能继续套原来的写法。

\textbf{相邻状态。} 规律是堆保存每组当前元素，同时维护当前最大值，最小头决定下一步推进哪一组。把这个特例继续拆开看：堆顶必须正好代表下一步最需要处理的候选；若堆里可能有过期对象，读取堆顶前要先清理。堆只解决动态极值，不自动证明被弹出或被丢弃的元素无用。

\textbf{状态复用。} 把每轮全量扫描改成堆顶查询；插入、弹出和过期清理共同维护动态候选集。本题落点是：规律是堆保存每组当前元素，同时维护当前最大值，最小头决定下一步推进哪一组。手算时只改被新输入影响的那一小部分，而不是回头重算完整候选。

\textbf{指针和字段。} 状态字段要能说清：堆元素字段、堆顶比较规则、候选是否过期、答案读取时机；如果需要懒删除，还要保存删除计数。手算时按这个协议跟踪：拿三组 [4, 10]、[0, 9]、[5, 18] 手算，当前头是 4、0、5，区间 [0, 5] 覆盖三组；弹出最小头 0 后推进该组。然后按协议复查：手算时记录堆顶、候选集合和是否有过期元素。每次使用堆顶前都要确认它仍然属于当前问题。

\textbf{代码落点。} 比较器方向先用小样例验证；每次使用堆顶前清理过期项；堆中保存下标时要能回查原数据。关键代码行必须能逐句翻译成“旧状态怎样变成新状态”，否则就只是模板。

\textbf{正确性。} 证明从候选覆盖开始：所有可能成为下一步答案的对象都在堆中，堆顶过期时不会被使用。

\textbf{边界实验。} 先用某个列表为空、重复值、区间长度相同的 tie-breaker、列表耗尽做反例检查。实验时覆盖某个列表为空、重复值、区间长度相同的 tie-breaker、列表耗尽，记录堆顶是否过期、比较器是否让正确候选在堆顶、K 或窗口大小为边界值时是否稳定。

\textbf{复杂度变化。} 暴力每轮扫描或排序动态候选，会把线性扫描或排序反复发生；堆把取极值、插入和删除边界控制在对数级。

\textbf{迁移追问。} 如果题目改成“这是多路归并和覆盖窗口的结合”，先判断原状态是否仍然覆盖新答案集合，再决定是微调边界、改 value 语义，还是换算法。

\noindent\textbf{LeetCode 703 数据流中的第 K 大元素。}

\textbf{题目说明。} LeetCode 703 数据流中的第 K 大元素 的题面入口要先还原成具体任务：数据不断到达，只关心当前第 K 大而非完整排序。

\textbf{样例入口。} 拿 k=3，初始 4, 5, 8, 2 手算，小顶堆保存 4, 5, 8，堆顶 4 是第 3 大；add 3 不进入堆，add 5 会替换 4。

\textbf{暴力基线。} 慢版本每轮扫描全部候选来找最大或最小，再更新候选集合。它的答案选择过程清楚，但动态集合每变化一次就重新找极值，代价太高。放到本题就是：先按“数据不断到达，只关心当前第 K 大而非完整排序”完整试慢版本；样例里已经能看到“规律是数据流只保留前 k 大，堆顶就是当前第 k 大”，所以优化前必须先能说清慢版本枚举了哪些候选。

\textbf{暴力过程。} 拿 k=3，初始 4, 5, 8, 2 手算，小顶堆保存 4, 5, 8，堆顶 4 是第 3 大；add 3 不进入堆，add 5 会替换 4。

\textbf{重复工作。} 题面模型：数据不断到达，只关心当前第 K 大而非完整排序。暴力版的慢点要落到本题：慢版本会围绕“数据不断到达，只关心当前第 K 大而非完整排序”反复扫描动态候选集找极值。样例规律是“规律是数据流只保留前 k 大，堆顶就是当前第 k 大”，说明每轮真正需要的是堆顶边界，而不是把候选全集重新排序。后面的优化只消掉这类重复工作，不能改变答案集合。

\textbf{优化条件。} 本题的关键观察是：维护大小为 K 的小顶堆，堆顶就是当前第 K 大元素。这句话必须能翻译成可维护状态；如果题目条件改变后观察不再成立，就不能继续套原来的写法。

\textbf{相邻状态。} 规律是数据流只保留前 k 大，堆顶就是当前第 k 大。把这个特例继续拆开看：堆顶必须正好代表下一步最需要处理的候选；若堆里可能有过期对象，读取堆顶前要先清理。堆只解决动态极值，不自动证明被弹出或被丢弃的元素无用。

\textbf{状态复用。} 把每轮全量扫描改成堆顶查询；插入、弹出和过期清理共同维护动态候选集。本题落点是：规律是数据流只保留前 k 大，堆顶就是当前第 k 大。手算时只改被新输入影响的那一小部分，而不是回头重算完整候选。

\textbf{指针和字段。} 状态字段要能说清：堆元素字段、堆顶比较规则、候选是否过期、答案读取时机；如果需要懒删除，还要保存删除计数。手算时按这个协议跟踪：拿 k=3，初始 4, 5, 8, 2 手算，小顶堆保存 4, 5, 8，堆顶 4 是第 3 大；add 3 不进入堆，add 5 会替换 4。然后按协议复查：手算时记录堆顶、候选集合和是否有过期元素。每次使用堆顶前都要确认它仍然属于当前问题。

\textbf{代码落点。} 比较器方向先用小样例验证；每次使用堆顶前清理过期项；堆中保存下标时要能回查原数据。关键代码行必须能逐句翻译成“旧状态怎样变成新状态”，否则就只是模板。

\textbf{正确性。} 证明从候选覆盖开始：所有可能成为下一步答案的对象都在堆中，堆顶过期时不会被使用。

\textbf{边界实验。} 先用初始元素少于 K、重复值、K 为一、连续 add做反例检查。实验时覆盖初始元素少于 K、重复值、K 为一、连续 add，记录堆顶是否过期、比较器是否让正确候选在堆顶、K 或窗口大小为边界值时是否稳定。

\textbf{复杂度变化。} 暴力每轮扫描或排序动态候选，会把线性扫描或排序反复发生；堆把取极值、插入和删除边界控制在对数级。

\textbf{迁移追问。} 如果题目改成“这题是 Top K 的在线版本，和一次性排序的适用场景不同”，先判断原状态是否仍然覆盖新答案集合，再决定是微调边界、改 value 语义，还是换算法。

\noindent\textbf{LeetCode 767 重构字符串。}

\textbf{题目说明。} LeetCode 767 重构字符串 的题面入口要先还原成具体任务：相邻不能相同，剩余次数最多的字符最难安排。

\textbf{样例入口。} 拿 s=aab 手算，a 剩 2 次、b 剩 1 次，每轮取两个最高频字符放入答案，可得到 ab 后再放回剩余 a，最后是 aba。拿 aaab 时 a 出现 3 次超过 (4+1)/2，不可能隔开。

\textbf{暴力基线。} 慢版本每轮扫描全部候选来找最大或最小，再更新候选集合。它的答案选择过程清楚，但动态集合每变化一次就重新找极值，代价太高。放到本题就是：先按“相邻不能相同，剩余次数最多的字符最难安排”完整试慢版本；样例里已经能看到“规律是最大频次先判不可能，构造时堆里每轮取两个不同最高频字符”，所以优化前必须先能说清慢版本枚举了哪些候选。

\textbf{暴力过程。} 拿 s=aab 手算，a 剩 2 次、b 剩 1 次，每轮取两个最高频字符放入答案，可得到 ab 后再放回剩余 a，最后是 aba。拿 aaab 时 a 出现 3 次超过 (4+1)/2，不可能隔开。

\textbf{重复工作。} 题面模型：相邻不能相同，剩余次数最多的字符最难安排。暴力版的慢点要落到本题：慢版本会围绕“相邻不能相同，剩余次数最多的字符最难安排”反复扫描动态候选集找极值。样例规律是“规律是最大频次先判不可能，构造时堆里每轮取两个不同最高频字符”，说明每轮真正需要的是堆顶边界，而不是把候选全集重新排序。后面的优化只消掉这类重复工作，不能改变答案集合。

\textbf{优化条件。} 本题的关键观察是：大顶堆每次取两个最高频不同字符放入答案，剩余频次再放回；先检查最大频次是否超过 (n+1)/2。这句话必须能翻译成可维护状态；如果题目条件改变后观察不再成立，就不能继续套原来的写法。

\textbf{相邻状态。} 规律是最大频次先判不可能，构造时堆里每轮取两个不同最高频字符。把这个特例继续拆开看：堆顶必须正好代表下一步最需要处理的候选；若堆里可能有过期对象，读取堆顶前要先清理。堆只解决动态极值，不自动证明被弹出或被丢弃的元素无用。

\textbf{状态复用。} 把每轮全量扫描改成堆顶查询；插入、弹出和过期清理共同维护动态候选集。本题落点是：规律是最大频次先判不可能，构造时堆里每轮取两个不同最高频字符。手算时只改被新输入影响的那一小部分，而不是回头重算完整候选。

\textbf{指针和字段。} 状态字段要能说清：堆元素字段、堆顶比较规则、候选是否过期、答案读取时机；如果需要懒删除，还要保存删除计数。手算时按这个协议跟踪：拿 s=aab 手算，a 剩 2 次、b 剩 1 次，每轮取两个最高频字符放入答案，可得到 ab 后再放回剩余 a，最后是 aba。拿 aaab 时 a 出现 3 次超过 (4+1)/2，不可能隔开。然后按协议复查：手算时记录堆顶、候选集合和是否有过期元素。每次使用堆顶前都要确认它仍然属于当前问题。

\textbf{代码落点。} 比较器方向先用小样例验证；每次使用堆顶前清理过期项；堆中保存下标时要能回查原数据。关键代码行必须能逐句翻译成“旧状态怎样变成新状态”，否则就只是模板。

\textbf{正确性。} 证明从候选覆盖开始：所有可能成为下一步答案的对象都在堆中，堆顶过期时不会被使用。

\textbf{边界实验。} 先用单字符、最高频过半、多个字符同频、最后只剩一个字符、结果不唯一做反例检查。实验时覆盖单字符、最高频过半、多个字符同频、最后只剩一个字符、结果不唯一，记录堆顶是否过期、比较器是否让正确候选在堆顶、K 或窗口大小为边界值时是否稳定。

\textbf{复杂度变化。} 暴力每轮扫描或排序动态候选，会把线性扫描或排序反复发生；堆把取极值、插入和删除边界控制在对数级。

\textbf{迁移追问。} 如果题目改成“若要求相同字符至少间隔 k，堆还要配合冷却队列”，先判断原状态是否仍然覆盖新答案集合，再决定是微调边界、改 value 语义，还是换算法。

\noindent\textbf{LeetCode 857 雇佣 K 名工人的最低成本。}

\textbf{题目说明。} LeetCode 857 雇佣 K 名工人的最低成本 的题面入口要先还原成具体任务：工资比例由当前队伍中最高工资质量比决定，总成本等于比例乘质量和。

\textbf{样例入口。} 拿 quality=[10, 20, 5]、wage=[70, 50, 30] 手算，按工资质量比选组时，当前最大比率决定组内所有人的单位工资；再用堆保留质量和最小的 k 人。

\textbf{暴力基线。} 慢版本每轮扫描全部候选来找最大或最小，再更新候选集合。它的答案选择过程清楚，但动态集合每变化一次就重新找极值，代价太高。放到本题就是：先按“工资比例由当前队伍中最高工资质量比决定，总成本等于比例乘质量和”完整试慢版本；样例里已经能看到“规律是排序枚举最高比率，堆优化质量总和”，所以优化前必须先能说清慢版本枚举了哪些候选。

\textbf{暴力过程。} 拿 quality=[10, 20, 5]、wage=[70, 50, 30] 手算，按工资质量比选组时，当前最大比率决定组内所有人的单位工资；再用堆保留质量和最小的 k 人。

\textbf{重复工作。} 题面模型：工资比例由当前队伍中最高工资质量比决定，总成本等于比例乘质量和。暴力版的慢点要落到本题：慢版本会围绕“工资比例由当前队伍中最高工资质量比决定，总成本等于比例乘质量和”反复扫描动态候选集找极值。样例规律是“规律是排序枚举最高比率，堆优化质量总和”，说明每轮真正需要的是堆顶边界，而不是把候选全集重新排序。后面的优化只消掉这类重复工作，不能改变答案集合。

\textbf{优化条件。} 本题的关键观察是：按比例升序扫描，用大顶堆维护当前 K 个工人的最小质量和。这句话必须能翻译成可维护状态；如果题目条件改变后观察不再成立，就不能继续套原来的写法。

\textbf{相邻状态。} 规律是排序枚举最高比率，堆优化质量总和。把这个特例继续拆开看：堆顶必须正好代表下一步最需要处理的候选；若堆里可能有过期对象，读取堆顶前要先清理。堆只解决动态极值，不自动证明被弹出或被丢弃的元素无用。

\textbf{状态复用。} 把每轮全量扫描改成堆顶查询；插入、弹出和过期清理共同维护动态候选集。本题落点是：规律是排序枚举最高比率，堆优化质量总和。手算时只改被新输入影响的那一小部分，而不是回头重算完整候选。

\textbf{指针和字段。} 状态字段要能说清：堆元素字段、堆顶比较规则、候选是否过期、答案读取时机；如果需要懒删除，还要保存删除计数。手算时按这个协议跟踪：拿 quality=[10, 20, 5]、wage=[70, 50, 30] 手算，按工资质量比选组时，当前最大比率决定组内所有人的单位工资；再用堆保留质量和最小的 k 人。然后按协议复查：手算时记录堆顶、候选集合和是否有过期元素。每次使用堆顶前都要确认它仍然属于当前问题。

\textbf{代码落点。} 比较器方向先用小样例验证；每次使用堆顶前清理过期项；堆中保存下标时要能回查原数据。关键代码行必须能逐句翻译成“旧状态怎样变成新状态”，否则就只是模板。

\textbf{正确性。} 证明从候选覆盖开始：所有可能成为下一步答案的对象都在堆中，堆顶过期时不会被使用。

\textbf{边界实验。} 先用K 为一、比例相同、浮点精度、质量和更新做反例检查。实验时覆盖K 为一、比例相同、浮点精度、质量和更新，记录堆顶是否过期、比较器是否让正确候选在堆顶、K 或窗口大小为边界值时是否稳定。

\textbf{复杂度变化。} 暴力每轮扫描或排序动态候选，会把线性扫描或排序反复发生；堆把取极值、插入和删除边界控制在对数级。

\textbf{迁移追问。} 如果题目改成“这是排序确定价格比例，堆维护可替换候选的组合题”，先判断原状态是否仍然覆盖新答案集合，再决定是微调边界、改 value 语义，还是换算法。

\noindent\textbf{LeetCode 871 最低加油次数。}

\textbf{题目说明。} LeetCode 871 最低加油次数 的题面入口要先还原成具体任务：行驶过程中，经过的加油站形成可反悔的燃料候选。

\textbf{样例入口。} 拿 target=100、startFuel=10 手算，车先到 10 号站，把 60 升油加入最大堆；之后油量不足时，从已经路过的站里取当前最大油量。

\textbf{暴力基线。} 慢版本每轮扫描全部候选来找最大或最小，再更新候选集合。它的答案选择过程清楚，但动态集合每变化一次就重新找极值，代价太高。放到本题就是：先按“行驶过程中，经过的加油站形成可反悔的燃料候选”完整试慢版本；样例里已经能看到“规律是能到的站先入堆，需要加油时再反悔选择最大燃料，保证次数最少”，所以优化前必须先能说清慢版本枚举了哪些候选。

\textbf{暴力过程。} 拿 target=100、startFuel=10 手算，车先到 10 号站，把 60 升油加入最大堆；之后油量不足时，从已经路过的站里取当前最大油量。

\textbf{重复工作。} 题面模型：行驶过程中，经过的加油站形成可反悔的燃料候选。暴力版的慢点要落到本题：慢版本会围绕“行驶过程中，经过的加油站形成可反悔的燃料候选”反复扫描动态候选集找极值。样例规律是“规律是能到的站先入堆，需要加油时再反悔选择最大燃料，保证次数最少”，说明每轮真正需要的是堆顶边界，而不是把候选全集重新排序。后面的优化只消掉这类重复工作，不能改变答案集合。

\textbf{优化条件。} 本题的关键观察是：按位置扫描，油量不足到达下一点时，从已路过站点中取最大燃料补充。这句话必须能翻译成可维护状态；如果题目条件改变后观察不再成立，就不能继续套原来的写法。

\textbf{相邻状态。} 规律是能到的站先入堆，需要加油时再反悔选择最大燃料，保证次数最少。把这个特例继续拆开看：堆顶必须正好代表下一步最需要处理的候选；若堆里可能有过期对象，读取堆顶前要先清理。堆只解决动态极值，不自动证明被弹出或被丢弃的元素无用。

\textbf{状态复用。} 把每轮全量扫描改成堆顶查询；插入、弹出和过期清理共同维护动态候选集。本题落点是：规律是能到的站先入堆，需要加油时再反悔选择最大燃料，保证次数最少。手算时只改被新输入影响的那一小部分，而不是回头重算完整候选。

\textbf{指针和字段。} 状态字段要能说清：堆元素字段、堆顶比较规则、候选是否过期、答案读取时机；如果需要懒删除，还要保存删除计数。手算时按这个协议跟踪：拿 target=100、startFuel=10 手算，车先到 10 号站，把 60 升油加入最大堆；之后油量不足时，从已经路过的站里取当前最大油量。然后按协议复查：手算时记录堆顶、候选集合和是否有过期元素。每次使用堆顶前都要确认它仍然属于当前问题。

\textbf{代码落点。} 比较器方向先用小样例验证；每次使用堆顶前清理过期项；堆中保存下标时要能回查原数据。关键代码行必须能逐句翻译成“旧状态怎样变成新状态”，否则就只是模板。

\textbf{正确性。} 证明从候选覆盖开始：所有可能成为下一步答案的对象都在堆中，堆顶过期时不会被使用。

\textbf{边界实验。} 先用起点油足够、无法到达第一个站、多个站同位置、目标作为终点事件做反例检查。实验时覆盖起点油足够、无法到达第一个站、多个站同位置、目标作为终点事件，记录堆顶是否过期、比较器是否让正确候选在堆顶、K 或窗口大小为边界值时是否稳定。

\textbf{复杂度变化。} 暴力每轮扫描或排序动态候选，会把线性扫描或排序反复发生；堆把取极值、插入和删除边界控制在对数级。

\textbf{迁移追问。} 如果题目改成“这是反悔贪心，堆里保存的是已经经过但尚未使用的资源”，先判断原状态是否仍然覆盖新答案集合，再决定是微调边界、改 value 语义，还是换算法。

\noindent\textbf{LeetCode 973 最接近原点的 K 个点。}

\textbf{题目说明。} LeetCode 973 最接近原点的 K 个点 的题面入口要先还原成具体任务：距离只需比较平方，避免开方误差和成本。

\textbf{样例入口。} 拿点 (1, 3) 和 (-2, 2) 手算，距离平方分别是 10 和 8，所以 k=1 时选 (-2, 2)。

\textbf{暴力基线。} 慢版本每轮扫描全部候选来找最大或最小，再更新候选集合。它的答案选择过程清楚，但动态集合每变化一次就重新找极值，代价太高。放到本题就是：先按“距离只需比较平方，避免开方误差和成本”完整试慢版本；样例里已经能看到“规律是比较距离平方避免开方，用大小为 k 的大顶堆保留当前最近点”，所以优化前必须先能说清慢版本枚举了哪些候选。

\textbf{暴力过程。} 拿点 (1, 3) 和 (-2, 2) 手算，距离平方分别是 10 和 8，所以 k=1 时选 (-2, 2)。

\textbf{重复工作。} 题面模型：距离只需比较平方，避免开方误差和成本。暴力版的慢点要落到本题：慢版本会围绕“距离只需比较平方，避免开方误差和成本”反复扫描动态候选集找极值。样例规律是“规律是比较距离平方避免开方，用大小为 k 的大顶堆保留当前最近点”，说明每轮真正需要的是堆顶边界，而不是把候选全集重新排序。后面的优化只消掉这类重复工作，不能改变答案集合。

\textbf{优化条件。} 本题的关键观察是：大小为 K 的大顶堆保存当前最近 K 个点中最远者。这句话必须能翻译成可维护状态；如果题目条件改变后观察不再成立，就不能继续套原来的写法。

\textbf{相邻状态。} 规律是比较距离平方避免开方，用大小为 k 的大顶堆保留当前最近点。把这个特例继续拆开看：堆顶必须正好代表下一步最需要处理的候选；若堆里可能有过期对象，读取堆顶前要先清理。堆只解决动态极值，不自动证明被弹出或被丢弃的元素无用。

\textbf{状态复用。} 把每轮全量扫描改成堆顶查询；插入、弹出和过期清理共同维护动态候选集。本题落点是：规律是比较距离平方避免开方，用大小为 k 的大顶堆保留当前最近点。手算时只改被新输入影响的那一小部分，而不是回头重算完整候选。

\textbf{指针和字段。} 状态字段要能说清：堆元素字段、堆顶比较规则、候选是否过期、答案读取时机；如果需要懒删除，还要保存删除计数。手算时按这个协议跟踪：拿点 (1, 3) 和 (-2, 2) 手算，距离平方分别是 10 和 8，所以 k=1 时选 (-2, 2)。然后按协议复查：手算时记录堆顶、候选集合和是否有过期元素。每次使用堆顶前都要确认它仍然属于当前问题。

\textbf{代码落点。} 比较器方向先用小样例验证；每次使用堆顶前清理过期项；堆中保存下标时要能回查原数据。关键代码行必须能逐句翻译成“旧状态怎样变成新状态”，否则就只是模板。

\textbf{正确性。} 证明从候选覆盖开始：所有可能成为下一步答案的对象都在堆中，堆顶过期时不会被使用。

\textbf{边界实验。} 先用K 为一、K 为全部、距离相同、坐标平方溢出做反例检查。实验时覆盖K 为一、K 为全部、距离相同、坐标平方溢出，记录堆顶是否过期、比较器是否让正确候选在堆顶、K 或窗口大小为边界值时是否稳定。

\textbf{复杂度变化。} 暴力每轮扫描或排序动态候选，会把线性扫描或排序反复发生；堆把取极值、插入和删除边界控制在对数级。

\textbf{迁移追问。} 如果题目改成“快速选择可平均线性，但堆更适合数据流”，先判断原状态是否仍然覆盖新答案集合，再决定是微调边界、改 value 语义，还是换算法。

\noindent\textbf{LeetCode 1046 最后一块石头的重量。}

\textbf{题目说明。} LeetCode 1046 最后一块石头的重量 的题面入口要先还原成具体任务：每次都要取当前最重的两块石头碰撞。

\textbf{样例入口。} 拿 [2, 7, 4, 1, 8, 1] 手算，每次取两块最大石头 8 和 7，撞剩 1 再放回。

\textbf{暴力基线。} 慢版本每轮扫描全部候选来找最大或最小，再更新候选集合。它的答案选择过程清楚，但动态集合每变化一次就重新找极值，代价太高。放到本题就是：先按“每次都要取当前最重的两块石头碰撞”完整试慢版本；样例里已经能看到“规律是动态最大值查询，用大顶堆保存当前石头重量，直到剩 0 或 1 块”，所以优化前必须先能说清慢版本枚举了哪些候选。

\textbf{暴力过程。} 拿 [2, 7, 4, 1, 8, 1] 手算，每次取两块最大石头 8 和 7，撞剩 1 再放回。

\textbf{重复工作。} 题面模型：每次都要取当前最重的两块石头碰撞。暴力版的慢点要落到本题：慢版本会围绕“每次都要取当前最重的两块石头碰撞”反复扫描动态候选集找极值。样例规律是“规律是动态最大值查询，用大顶堆保存当前石头重量，直到剩 0 或 1 块”，说明每轮真正需要的是堆顶边界，而不是把候选全集重新排序。后面的优化只消掉这类重复工作，不能改变答案集合。

\textbf{优化条件。} 本题的关键观察是：大顶堆保存所有重量，弹出两个最大值，若不同则把差值放回。这句话必须能翻译成可维护状态；如果题目条件改变后观察不再成立，就不能继续套原来的写法。

\textbf{相邻状态。} 规律是动态最大值查询，用大顶堆保存当前石头重量，直到剩 0 或 1 块。把这个特例继续拆开看：堆顶必须正好代表下一步最需要处理的候选；若堆里可能有过期对象，读取堆顶前要先清理。堆只解决动态极值，不自动证明被弹出或被丢弃的元素无用。

\textbf{状态复用。} 把每轮全量扫描改成堆顶查询；插入、弹出和过期清理共同维护动态候选集。本题落点是：规律是动态最大值查询，用大顶堆保存当前石头重量，直到剩 0 或 1 块。手算时只改被新输入影响的那一小部分，而不是回头重算完整候选。

\textbf{指针和字段。} 状态字段要能说清：堆元素字段、堆顶比较规则、候选是否过期、答案读取时机；如果需要懒删除，还要保存删除计数。手算时按这个协议跟踪：拿 [2, 7, 4, 1, 8, 1] 手算，每次取两块最大石头 8 和 7，撞剩 1 再放回。然后按协议复查：手算时记录堆顶、候选集合和是否有过期元素。每次使用堆顶前都要确认它仍然属于当前问题。

\textbf{代码落点。} 比较器方向先用小样例验证；每次使用堆顶前清理过期项；堆中保存下标时要能回查原数据。关键代码行必须能逐句翻译成“旧状态怎样变成新状态”，否则就只是模板。

\textbf{正确性。} 证明从候选覆盖开始：所有可能成为下一步答案的对象都在堆中，堆顶过期时不会被使用。

\textbf{边界实验。} 先用没有石头、一块石头、两块相等、多次产生相同重量做反例检查。实验时覆盖没有石头、一块石头、两块相等、多次产生相同重量，记录堆顶是否过期、比较器是否让正确候选在堆顶、K 或窗口大小为边界值时是否稳定。

\textbf{复杂度变化。} 暴力每轮扫描或排序动态候选，会把线性扫描或排序反复发生；堆把取极值、插入和删除边界控制在对数级。

\textbf{迁移追问。} 如果题目改成“最后一块石头 II 则是分组差值最小化，模型转为 0-1 背包”，先判断原状态是否仍然覆盖新答案集合，再决定是微调边界、改 value 语义，还是换算法。

\noindent\textbf{LeetCode 1353 最多可以参加的会议数目。}

\textbf{题目说明。} LeetCode 1353 最多可以参加的会议数目 的题面入口要先还原成具体任务：每天只能参加一个正在进行的会议，应优先选择最早结束的会议。

\textbf{样例入口。} 拿 events=[[1, 2], [2, 3], [3, 4]] 手算，第 1 天把开始的会议入堆并参加最早结束的；第 2 天再清理过期并参加一个。

\textbf{暴力基线。} 慢版本每轮扫描全部候选来找最大或最小，再更新候选集合。它的答案选择过程清楚，但动态集合每变化一次就重新找极值，代价太高。放到本题就是：先按“每天只能参加一个正在进行的会议，应优先选择最早结束的会议”完整试慢版本；样例里已经能看到“规律是按天扫描，用小顶堆保存当前可参加会议的结束日，每天选最早结束者”，所以优化前必须先能说清慢版本枚举了哪些候选。

\textbf{暴力过程。} 拿 events=[[1, 2], [2, 3], [3, 4]] 手算，第 1 天把开始的会议入堆并参加最早结束的；第 2 天再清理过期并参加一个。

\textbf{重复工作。} 题面模型：每天只能参加一个正在进行的会议，应优先选择最早结束的会议。暴力版的慢点要落到本题：慢版本会围绕“每天只能参加一个正在进行的会议，应优先选择最早结束的会议”反复扫描动态候选集找极值。样例规律是“规律是按天扫描，用小顶堆保存当前可参加会议的结束日，每天选最早结束者”，说明每轮真正需要的是堆顶边界，而不是把候选全集重新排序。后面的优化只消掉这类重复工作，不能改变答案集合。

\textbf{优化条件。} 本题的关键观察是：按开始日排序扫描日期，把当天可参加会议的结束日入小顶堆，弹出最早结束者。这句话必须能翻译成可维护状态；如果题目条件改变后观察不再成立，就不能继续套原来的写法。

\textbf{相邻状态。} 规律是按天扫描，用小顶堆保存当前可参加会议的结束日，每天选最早结束者。把这个特例继续拆开看：堆顶必须正好代表下一步最需要处理的候选；若堆里可能有过期对象，读取堆顶前要先清理。堆只解决动态极值，不自动证明被弹出或被丢弃的元素无用。

\textbf{状态复用。} 把每轮全量扫描改成堆顶查询；插入、弹出和过期清理共同维护动态候选集。本题落点是：规律是按天扫描，用小顶堆保存当前可参加会议的结束日，每天选最早结束者。手算时只改被新输入影响的那一小部分，而不是回头重算完整候选。

\textbf{指针和字段。} 状态字段要能说清：堆元素字段、堆顶比较规则、候选是否过期、答案读取时机；如果需要懒删除，还要保存删除计数。手算时按这个协议跟踪：拿 events=[[1, 2], [2, 3], [3, 4]] 手算，第 1 天把开始的会议入堆并参加最早结束的；第 2 天再清理过期并参加一个。然后按协议复查：手算时记录堆顶、候选集合和是否有过期元素。每次使用堆顶前都要确认它仍然属于当前问题。

\textbf{代码落点。} 比较器方向先用小样例验证；每次使用堆顶前清理过期项；堆中保存下标时要能回查原数据。关键代码行必须能逐句翻译成“旧状态怎样变成新状态”，否则就只是模板。

\textbf{正确性。} 证明从候选覆盖开始：所有可能成为下一步答案的对象都在堆中，堆顶过期时不会被使用。

\textbf{边界实验。} 先用同一天多个会议、过期会议清理、空闲日期跳跃、结束日相同做反例检查。实验时覆盖同一天多个会议、过期会议清理、空闲日期跳跃、结束日相同，记录堆顶是否过期、比较器是否让正确候选在堆顶、K 或窗口大小为边界值时是否稳定。

\textbf{复杂度变化。} 暴力每轮扫描或排序动态候选，会把线性扫描或排序反复发生；堆把取极值、插入和删除边界控制在对数级。

\textbf{迁移追问。} 如果题目改成“这是扫描线和堆结合的调度题”，先判断原状态是否仍然覆盖新答案集合，再决定是微调边界、改 value 语义，还是换算法。

\noindent\textbf{LeetCode 1675 数组的最小偏移量。}

\textbf{题目说明。} LeetCode 1675 数组的最小偏移量 的题面入口要先还原成具体任务：奇数可以先乘二，偶数可以不断除二，目标是缩小当前最大最小差。

\textbf{样例入口。} 拿 [1, 2, 3, 4] 手算，奇数可先翻倍成偶数，之后只能不断把当前最大偶数减半；每次记录最大值和当前最小值差。

\textbf{暴力基线。} 慢版本每轮扫描全部候选来找最大或最小，再更新候选集合。它的答案选择过程清楚，但动态集合每变化一次就重新找极值，代价太高。放到本题就是：先按“奇数可以先乘二，偶数可以不断除二，目标是缩小当前最大最小差”完整试慢版本；样例里已经能看到“规律是把所有数变到可下降的最高状态，再用大顶堆逐步降低最大值”，所以优化前必须先能说清慢版本枚举了哪些候选。

\textbf{暴力过程。} 拿 [1, 2, 3, 4] 手算，奇数可先翻倍成偶数，之后只能不断把当前最大偶数减半；每次记录最大值和当前最小值差。

\textbf{重复工作。} 题面模型：奇数可以先乘二，偶数可以不断除二，目标是缩小当前最大最小差。暴力版的慢点要落到本题：慢版本会围绕“奇数可以先乘二，偶数可以不断除二，目标是缩小当前最大最小差”反复扫描动态候选集找极值。样例规律是“规律是把所有数变到可下降的最高状态，再用大顶堆逐步降低最大值”，说明每轮真正需要的是堆顶边界，而不是把候选全集重新排序。后面的优化只消掉这类重复工作，不能改变答案集合。

\textbf{优化条件。} 本题的关键观察是：把所有数规范到可降低的偶数形态，用大顶堆维护当前最大值，同时记录当前最小值。这句话必须能翻译成可维护状态；如果题目条件改变后观察不再成立，就不能继续套原来的写法。

\textbf{相邻状态。} 规律是把所有数变到可下降的最高状态，再用大顶堆逐步降低最大值。把这个特例继续拆开看：堆顶必须正好代表下一步最需要处理的候选；若堆里可能有过期对象，读取堆顶前要先清理。堆只解决动态极值，不自动证明被弹出或被丢弃的元素无用。

\textbf{状态复用。} 把每轮全量扫描改成堆顶查询；插入、弹出和过期清理共同维护动态候选集。本题落点是：规律是把所有数变到可下降的最高状态，再用大顶堆逐步降低最大值。手算时只改被新输入影响的那一小部分，而不是回头重算完整候选。

\textbf{指针和字段。} 状态字段要能说清：堆元素字段、堆顶比较规则、候选是否过期、答案读取时机；如果需要懒删除，还要保存删除计数。手算时按这个协议跟踪：拿 [1, 2, 3, 4] 手算，奇数可先翻倍成偶数，之后只能不断把当前最大偶数减半；每次记录最大值和当前最小值差。然后按协议复查：手算时记录堆顶、候选集合和是否有过期元素。每次使用堆顶前都要确认它仍然属于当前问题。

\textbf{代码落点。} 比较器方向先用小样例验证；每次使用堆顶前清理过期项；堆中保存下标时要能回查原数据。关键代码行必须能逐句翻译成“旧状态怎样变成新状态”，否则就只是模板。

\textbf{正确性。} 证明从候选覆盖开始：所有可能成为下一步答案的对象都在堆中，堆顶过期时不会被使用。

\textbf{边界实验。} 先用全奇数、全偶数、最大值变奇数停止、数值范围做反例检查。实验时覆盖全奇数、全偶数、最大值变奇数停止、数值范围，记录堆顶是否过期、比较器是否让正确候选在堆顶、K 或窗口大小为边界值时是否稳定。

\textbf{复杂度变化。} 暴力每轮扫描或排序动态候选，会把线性扫描或排序反复发生；堆把取极值、插入和删除边界控制在对数级。

\textbf{迁移追问。} 如果题目改成“这是堆维护动态极值和状态空间规范化的结合”，先判断原状态是否仍然覆盖新答案集合，再决定是微调边界、改 value 语义，还是换算法。

\noindent\textbf{LeetCode 1851 包含每个查询的最小区间。}

\textbf{题目说明。} LeetCode 1851 包含每个查询的最小区间 的题面入口要先还原成具体任务：每个查询点需要找覆盖它的最短区间，区间和查询都可离线排序。

\textbf{样例入口。} 拿 intervals=[[1, 4], [2, 4], [3, 6]]、queries=[2, 3, 5] 手算，查询 2 时可选 [1, 4] 和 [2, 4]，较短的是 [2, 4]。

\textbf{暴力基线。} 慢版本每轮扫描全部候选来找最大或最小，再更新候选集合。它的答案选择过程清楚，但动态集合每变化一次就重新找极值，代价太高。放到本题就是：先按“每个查询点需要找覆盖它的最短区间，区间和查询都可离线排序”完整试慢版本；样例里已经能看到“规律是离线按查询升序加入左端不大于查询的区间，堆顶清理右端已过期区间”，所以优化前必须先能说清慢版本枚举了哪些候选。

\textbf{暴力过程。} 拿 intervals=[[1, 4], [2, 4], [3, 6]]、queries=[2, 3, 5] 手算，查询 2 时可选 [1, 4] 和 [2, 4]，较短的是 [2, 4]。

\textbf{重复工作。} 题面模型：每个查询点需要找覆盖它的最短区间，区间和查询都可离线排序。暴力版的慢点要落到本题：慢版本会围绕“每个查询点需要找覆盖它的最短区间，区间和查询都可离线排序”反复扫描动态候选集找极值。样例规律是“规律是离线按查询升序加入左端不大于查询的区间，堆顶清理右端已过期区间”，说明每轮真正需要的是堆顶边界，而不是把候选全集重新排序。后面的优化只消掉这类重复工作，不能改变答案集合。

\textbf{优化条件。} 本题的关键观察是：按查询值升序扫描，把左端不大于查询的区间入堆，堆顶过期区间弹出。这句话必须能翻译成可维护状态；如果题目条件改变后观察不再成立，就不能继续套原来的写法。

\textbf{相邻状态。} 规律是离线按查询升序加入左端不大于查询的区间，堆顶清理右端已过期区间。把这个特例继续拆开看：堆顶必须正好代表下一步最需要处理的候选；若堆里可能有过期对象，读取堆顶前要先清理。堆只解决动态极值，不自动证明被弹出或被丢弃的元素无用。

\textbf{状态复用。} 把每轮全量扫描改成堆顶查询；插入、弹出和过期清理共同维护动态候选集。本题落点是：规律是离线按查询升序加入左端不大于查询的区间，堆顶清理右端已过期区间。手算时只改被新输入影响的那一小部分，而不是回头重算完整候选。

\textbf{指针和字段。} 状态字段要能说清：堆元素字段、堆顶比较规则、候选是否过期、答案读取时机；如果需要懒删除，还要保存删除计数。手算时按这个协议跟踪：拿 intervals=[[1, 4], [2, 4], [3, 6]]、queries=[2, 3, 5] 手算，查询 2 时可选 [1, 4] 和 [2, 4]，较短的是 [2, 4]。然后按协议复查：手算时记录堆顶、候选集合和是否有过期元素。每次使用堆顶前都要确认它仍然属于当前问题。

\textbf{代码落点。} 比较器方向先用小样例验证；每次使用堆顶前清理过期项；堆中保存下标时要能回查原数据。关键代码行必须能逐句翻译成“旧状态怎样变成新状态”，否则就只是模板。

\textbf{正确性。} 证明从候选覆盖开始：所有可能成为下一步答案的对象都在堆中，堆顶过期时不会被使用。

\textbf{边界实验。} 先用没有覆盖区间、区间长度相同、查询重复、右端过期做反例检查。实验时覆盖没有覆盖区间、区间长度相同、查询重复、右端过期，记录堆顶是否过期、比较器是否让正确候选在堆顶、K 或窗口大小为边界值时是否稳定。

\textbf{复杂度变化。} 暴力每轮扫描或排序动态候选，会把线性扫描或排序反复发生；堆把取极值、插入和删除边界控制在对数级。

\textbf{迁移追问。} 如果题目改成“这是离线扫描线、堆和区间覆盖的组合题”，先判断原状态是否仍然覆盖新答案集合，再决定是微调边界、改 value 语义，还是换算法。

\subsection{实验复盘}

追问通常会问为什么不完整排序、堆顶是否可能过期、比较器方向是否写反。请准备说明堆里元素的业务含义，而不是只说 priority\_queue。

复盘本节时先从 LeetCode 23 合并 K 个升序链表、LeetCode 215 数组中的第 K 个最大元素、LeetCode 295 数据流的中位数 开始：每道题都先手写一个能覆盖全部候选的暴力版，再指出优化结构具体保存了哪一段历史信息，最后用相邻案例比较为什么不能互换解法。
